\booksection{第一章附：鸿蒙操作系统专题}

\chapterintro{鸿蒙（HarmonyOS）是华为自主研发的操作系统，从2019年发布至今，已经走过了几个大版本。它不只是“又一个安卓替代品”，而是有着独特设计理念的新物种。让我们一起深入了解这个正在崛起的操作系统。}

\booksubsection{一、鸿蒙是什么？为什么它与众不同？}

鸿蒙的核心理念是“万物互联”。不同于传统操作系统主要服务于单一设备，鸿蒙从设计之初就考虑了多设备协同。

\textbf{分布式架构}是鸿蒙的灵魂。简单说，它可以把多个设备虚拟成一个“超级终端”。你的手机、平板、电脑、手表、电视，不再是孤立的个体，而是一个整体的“器官”，可以互相调用对方的能力。

举个例子：你在用手机看视频，想在大屏上看，传统方式是投屏；而在鸿蒙上，你可以直接把视频“流转”到电视上继续看，手机变成遥控器。这不是简单的投屏，而是应用本身在跨设备运行。

\booksubsection{二、鸿蒙4.3：基于AOSP的成熟版本}

鸿蒙4.3是当前华为手机上广泛部署的版本。需要说明的是，鸿蒙4.3兼容安卓开源项目（AOSP）的代码框架，因此它能够完整兼容安卓应用——你在华为应用市场下载的大部分App，本质上仍然是安卓应用，通过鸿蒙的兼容层运行。但鸿蒙4.3并非简单套壳安卓，它在AOSP基础上加入了大量华为自研的分布式能力、方舟引擎等技术。

\subsubsection*{核心特性：}

\begin{itemize}
  \item \textbf{安卓应用兼容}：由于基于AOSP开发，鸿蒙4.3可以正常运行几乎所有安卓应用，用户几乎感受不到差异。
  \item \textbf{方舟引擎}：华为自研的编译器和运行时，让应用运行更流畅、更省电。相比早期版本，流畅度提升明显。
  \item \textbf{超级终端}：在控制中心，可以直观地看到周围的其他鸿蒙设备，拖拽即可连接。手机调用电脑的键盘鼠标、平板调用手机的摄像头，都很方便。
  \item \textbf{服务卡片}：类似iOS的小组件，但更灵活。可以上滑App图标展开卡片，显示关键信息，不用打开App。
  \item \textbf{纯净模式}：只允许从官方应用市场安装应用，减少恶意软件风险。
  \item \textbf{隐私中心}：集中管理应用权限，查看应用行为记录。
\end{itemize}

\subsubsection*{用户界面特点：}

\begin{itemize}
  \item \textbf{杂志锁屏}：锁屏界面显示精美图片和文章，划走即进入桌面。
  \item \textbf{万能卡片}：桌面图标可以展开成卡片，显示更多信息。
  \item \textbf{超级中转站}：一个临时存储区，可以把文字、图片、文件暂存，然后拖到任意App中使用。跨应用分享变得无比丝滑。
  \item \textbf{智慧搜索}：全局搜索，可以搜应用、文件、设置、网页。
\end{itemize}

\booksubsection{三、鸿蒙NEXT（鸿蒙6）：华为自主操作系统的真正起点}

鸿蒙NEXT（HarmonyOS NEXT 5.0）与鸿蒙4.3有着根本性的区别——它采用华为自研的鸿蒙微内核，不再依赖Linux内核和AOSP代码，也\textbf{不再直接支持安卓应用}。这是鸿蒙从“兼容安卓的过渡形态”正式走向“完全独立的操作系统”的关键一步。

\subsubsection*{重大变革：}

\begin{itemize}
  \item \textbf{完全自主内核}：鸿蒙NEXT使用华为自研的内核，不再包含任何AOSP代码。系统更精简、更安全、更高效。
  \item \textbf{不再支持安卓应用}：这是最大的变化。鸿蒙NEXT只能运行原生鸿蒙应用，无法安装和运行APK文件。升级前请确认你常用的App已经有鸿蒙原生版本。
  \item \textbf{原生应用生态}：华为正在全力推动开发者将应用迁移到鸿蒙原生平台。目前主流App如微信、支付宝、淘宝等已推出鸿蒙原生版本，但部分小众应用可能尚未适配。
  \item \textbf{星盾安全架构}：全新的安全体系，从硬件到软件的全链路保护。支持更细粒度的权限控制，比如"模糊定位"——应用只能知道你大概在哪，不知道精确位置。
  \item \textbf{AI深度融合}：系统级AI能力，不只是某个App有AI。比如AI字幕可以实时翻译视频，AI抠图可以智能识别主体。
  \item \textbf{跨设备无缝流转}：更进一步的分布式能力。应用状态可以在设备间无缝切换，你在手机上编辑的文档，走到电脑前自动在电脑上打开。
  \item \textbf{元服务}：无需安装即可使用的轻量级服务。比如扫二维码直接打开某个功能，不用先下载App。
\end{itemize}

\begin{tcbwarn}{升级鸿蒙NEXT前必读}
升级到鸿蒙NEXT后，你的手机将无法运行安卓应用。请务必在升级前确认：(1)你常用的App已有鸿蒙原生版本；(2)你的工作或生活中没有必须依赖安卓应用的场景。升级后无法回退到鸿蒙4.3。

\end{tcbwarn}

\subsubsection*{操作逻辑变化：}

\begin{itemize}
  \item \textbf{手势导航}：左右边缘滑动返回，底部上滑回桌面，上滑停顿进多任务。
  \item \textbf{智慧多窗}：支持分屏、悬浮窗、平行视界等多种多任务形态。
  \item \textbf{畅连通话}：华为的视频通话服务，支持屏幕共享、美颜、背景虚化。
\end{itemize}

\booksubsection{四、鸿蒙生态：设备与应用}

鸿蒙不只是手机系统，而是一个覆盖多种设备的生态系统。

\subsubsection*{支持的设备类型：}

\begin{itemize}
  \item \textbf{手机与平板}：核心设备，功能最完整。
  \item \textbf{智能手表与手环}：独立运行或与手机协同。
  \item \textbf{智慧屏（电视）}：家庭娱乐中心。
  \item \textbf{智能音箱}：语音交互入口。
  \item \textbf{汽车}：鸿蒙座舱，智能驾驶舱体验。
  \item \textbf{智能家居}：冰箱、空调、洗衣机等。
\end{itemize}

\subsubsection*{应用生态现状：}

鸿蒙应用市场已有数百万应用，主流App基本都有。华为也在积极推动开发者开发原生鸿蒙应用。截至2026年中，微信、支付宝、淘宝、抖音、美团、京东等头部应用均已推出鸿蒙原生版本，日常使用基本无障碍。部分小众应用和行业应用可能仍在适配中，升级前建议确认你常用的App是否已支持。

\booksubsection{五、鸿蒙实用技巧}

\begin{itemize}
  \item \textbf{超级终端}：从屏幕右上角下滑打开控制中心，点击"超级终端"，可以看到周围的其他设备，拖拽连接。
  \item \textbf{超级中转站}：长按文字/图片/文件，选择"复制到中转站"，然后在其他App中长按粘贴，或拖拽出来。
  \item \textbf{服务卡片}：上滑App图标展开卡片，长按卡片可以调整大小或添加到桌面。
  \item \textbf{智慧多窗}：从屏幕左右边缘向内滑动并停顿，打开侧边栏，点击应用以悬浮窗打开。
  \item \textbf{隐私保护}：设置 > 隐私，可以查看应用权限使用记录，设置模糊定位等。
\end{itemize}

\begin{tcbwarn}{注意事项}
鸿蒙系统在不断更新迭代，不同版本可能有些差异。本文描述基于鸿蒙4.3和6的主要特性，具体功能请以你的手机实际版本为准。建议保持系统更新，以获得最新功能和安全补丁。

\end{tcbwarn}

\begin{tcbpractice}{现在就试：体验鸿蒙特色功能}
\begin{enumerate}
  \item 打开控制中心，找到"超级终端"，看看周围有没有其他鸿蒙设备。
  \item 试试"超级中转站"：复制一段文字到中转站，然后在备忘录中拖出来。
  \item 上滑几个常用App图标，看看它们的服务卡片显示什么内容。
  \item 去设置 > 隐私，查看最近哪些应用使用了敏感权限。
\end{enumerate}

\end{tcbpractice}
